\documentclass[../../main.tex]{subfiles}
\graphicspath{{\subfix{../../image/}}}

\begin{document}

\section{古典概型}\label{section:ch01-古典概型}

\subsection*{一、模型与计算公式}

在讨论一般随机现象之前，我们先讨论一类最简单的随机现象。这种随机现象具有下列两个特征：

（1）在试验中它的全部可能结果只有有限个，譬如为 $n$ 个，记为 $\omega_1,\omega_2,\cdots,\omega_n$，而且这些事件是两两互不相容的；

（2）事件 $\omega_1,\omega_2,\cdots,\omega_n$ 的发生或出现是等可能的，即它们发生的概率都一样。

这类随机现象在概率论发展初期即被注意，许多最初的概率论结果也是对它作出的，一般把这类随机现象的数学模型称为\textbf{古典概型}。古典概型在概率论中占有相当重要的地位，一方面，由于它简单，对它的讨论有助于直观地理解概率论的许多基本概念，因此，我们常从讨论古典概型开始引入新的概念；另一方面，古典概型概率的计算在产品质量抽样检查等实际问题以及理论物理的研究中都有重要应用。

显然，古典概型是有限样本空间的一种特例。可以选 $\Omega=\{\omega_1,\omega_2,\cdots,\omega_n\}$ 作为样本空间，而且此时应有
\begin{align}
P(\omega_1)=P(\omega_2)=\cdots=P(\omega_n)=\frac{1}{n}\label{align:ch01-等可能}
\end{align}

对于任何事件 $A$，它总可以表示为样本点之和，例如 $A=\omega_{i_1}+\omega_{i_2}+\cdots+\omega_{i_m}$，因此由事件概率的定义
\begin{align}
P(A)&=P(\omega_{i_1})+P(\omega_{i_2})+\cdots+P(\omega_{i_m})\notag\\
&=\frac{1}{n}+\frac{1}{n}+\cdots+\frac{1}{n}=\frac{m}{n}\label{align:ch01-古典概率}
\end{align}

所以在古典概型中，事件 $A$ 的概率是一个分数，其分母是样本点的总数 $n$，而分子是事件 $A$ 中所包含的样本点的个数 $m$，由于 $\omega_{i_1},\omega_{i_2},\cdots,\omega_{i_m}$ 的出现必导致 $A$ 的出现，即它们的出现对 $A$ 的出现“有利”，因此习惯上常称 $\omega_{i_1},\omega_{i_2},\cdots,\omega_{i_m}$ 是 $A$ 的“有利场合”，这样，
\begin{align}
P(A)=\frac{m}{n}=\frac{A\text{ 的有利场合的数目}}{\text{样本点总数}}\label{align:ch01-古典概率定义}
\end{align}

法国数学家 Laplace（$1749$—$1827$）在 $1812$ 年把上式作为概率的一般定义。现在通常称它为概率的\textbf{古典定义}，因为它只适用于古典概型场合。

古典概型有着多方面应用，产品抽样检查就是其中之一。

产品抽样检查的技术，在各个生产部门中被广泛采用。许多大工厂产量很高，每天生产的产品数以万计，对这些产品的质量如果要进行全面的逐件检验通常是不可能的或是不经济的；另外，在有些情况下，产品的检验方法带有破坏性（如电灯泡寿命检验和棉纱强度试验），这样，最适宜的检验方法是采用抽样检查，即从产品中随机地抽出若干件来检验，并根据检验结果来判断整批产品的质量。

关于产品的质量，可以有多种多样的衡量标准，例如可能要考虑产品的某种形状或尺寸，或把产品分成若干等级，我们先考虑最简单的情形，即把产品分成合格品（好品）与次品（废品）两个类型的场合。

假如产品的好坏从外形上看不出来，而且我们又是随机抽样，那么任何一件产品被抽到的可能性都一样，这正是古典概型。

有一个口袋，内装 $a$ 只黑球，$b$ 只白球，它们除颜色不同外，外形完全一样（以后若非特别声明，均作此假定）。这样一来，当我们从袋子中任意摸出一球时，这 $a+b$ 只球中的任意一只被摸到的可能性都一样。

若把黑球作为废品，白球作为好品，则这个摸球模型就可以描述产品抽样。假如产品分为更多等级，例如一等品，二等品，三等品，等外品等等，则可用装有多种颜色的球的口袋的摸球模型来描述。

这种模型化的方法能使问题更清楚，更容易看出其随机性本质而不致被个别情况下的具体属性所蒙蔽。不仅如此，这种抽象化的模型带有普遍性，它还可以描述许多别的具体问题，从而有着多方面应用。例如种水稻地块的调查，某电视节目收视率的调查，某种疾病的抽查等都能用这个模型。

事实上，古典概型的大部分问题都能形象化地用摸球模型来描述。以后我们经常研究摸球模型，意义即在于此。

前节例~\ref{example:ch01-摸球} 及其有关概率的计算是古典概型的一个例子，但并不是所有古典概型的事件的概率计算都这么容易。事实上，古典概型中许多概率的计算相当困难而富有技巧。计算的要点是给定样本空间和样本点，并计算它的总数，而后再计算有利场合的数目。

在这些计算中，经常要用到一些排列与组合公式。

\subsection*{二、基本的组合分析公式}

\subsubsection*{1. 两条原理}

全部组合分析公式的推导基于下列两条原理：

\textbf{乘法原理}~~若进行 $A_1$ 过程有 $n_1$ 种方法，进行 $A_2$ 过程有 $n_2$ 种方法，则进行 $A_1$ 过程后再接着进行 $A_2$ 过程共有 $n_1\times n_2$ 种方法（图~\ref{figure:ch01-乘法原理}）。

\textbf{加法原理}~~若进行 $A_1$ 过程有 $n_1$ 种方法，进行 $A_2$ 过程有 $n_2$ 种方法，假定 $A_1$ 过程与 $A_2$ 过程是并行的，则进行过程 $A_1$ 或过程 $A_2$ 的方法共有 $n_1+n_2$ 种（图~\ref{figure:ch01-加法原理}）。

\begin{figure}[H]
\centering
\begin{minipage}{0.45\textwidth}
\centering
\includegraphics[width=\textwidth]{ch01-mult.png}
\caption{乘法原理}\label{figure:ch01-乘法原理}
\end{minipage}
\hfill
\begin{minipage}{0.45\textwidth}
\centering
\includegraphics[width=\textwidth]{ch01-add.png}
\caption{加法原理}\label{figure:ch01-加法原理}
\end{minipage}
\end{figure}

显然这二条原理可以拓广到多个过程的场合。

\subsubsection*{2. 排列}

从包含有 $n$ 个不同的元素的总体中取出 $r$ 个来进行排列，这时既要考虑到取出的元素也要顾及其取出顺序。

这种排列可分为两类：第一类是有放回地选取，这时每次选取都是在全体元素中进行，同一元素可被重复选中；另一类是不放回选取，这时一个元素一旦被取出便立刻从总体中除去，因此每个元素至多被选中一次，在后一类情况，必有 $r\leqslant n$。

（1）在有放回选取中，从 $n$ 个不同的元素中取出 $r$ 个元素进行排列，这种排列称为\textbf{有重复的排列}，其总数共有 $n^r$ 种。

（2）在不放回选取中，从 $n$ 个不同的元素中取出 $r$ 个元素进行排列，其总数为
\begin{align*}
A_n^r=n(n-1)(n-2)\cdots(n-r+1)
\end{align*}
这种排列称为\textbf{选排列}。特别当 $r=n$ 时，称为\textbf{全排列}。

（3）$n$ 个不同的元素的全排列数为
\begin{align*}
P_n=n(n-1)\cdots 3\cdot 2\cdot 1=n!
\end{align*}

\subsubsection*{3. 组合}

（1）从 $n$ 个不同的元素中取出 $r$ 个元素而不考虑其顺序，称为\textbf{组合}，其总数为
\begin{align*}
C_n^r=\binom{n}{r}=\frac{A_n^r}{r!}=\frac{n(n-1)\cdots(n-r+1)}{r!}=\frac{n!}{r!(n-r)!}
\end{align*}
这里 $\binom{n}{r}$ 称为\textbf{二项系数}，是下列二项展开式的系数：
\begin{align*}
(a+b)^n=\sum_{r=0}^{n}\binom{n}{r}a^r b^{n-r}.
\end{align*}

（2）若 $r_1+r_2+\cdots+r_k=n$，把 $n$ 个不同的元素分成 $k$ 个部分，第一部分 $r_1$ 个，第二部分 $r_2$ 个，$\cdots\cdots$，第 $k$ 部分 $r_k$ 个，则不同的分法有
\begin{align}
\frac{n!}{r_1!\,r_2!\,\cdots r_k!}\label{align:ch01-多项式系数}
\end{align}
种，上式中的数称为\textbf{多项式系数}，因为它是 $(x_1+x_2+\cdots+x_k)^n$ 展开式中 $x_1^{r_1}x_2^{r_2}\cdots x_k^{r_k}$ 的系数，当 $k=2$ 时，即为二项系数。

（3）若 $n$ 个元素中有 $n_1$ 个带足标“$1$”，$n_2$ 个带足标“$2$”，$\cdots\cdots$，$n_k$ 个带足标“$k$”，且 $n_1+n_2+\cdots+n_k=n$，从这 $n$ 个元素中取出 $r$ 个，使得带有足标“$i$”的元素有 $r_i$ 个（$1\leqslant i\leqslant k$），而 $r_1+r_2+\cdots+r_k=r$，这时不同取法的总数为
\begin{align}
\binom{n_1}{r_1}\binom{n_2}{r_2}\cdots\binom{n_k}{r_k}\label{align:ch01-分类组合}
\end{align}

（4）从 $n$ 个不同的元素中有重复地取 $r$ 个，不计顺序，则不同的取法有
\begin{align*}
\binom{n+r-1}{r}
\end{align*}
种，这个数称为\textbf{有重复组合数}。

\subsubsection*{4. 关于二项系数的一些公式}

在二项系数的定义式中，若约定 $0!=1$，则对一切 $0\leqslant k\leqslant n$，成立
\begin{align*}
\binom{n}{k}=\binom{n}{n-k}
\end{align*}
此外，对正整数 $n$ 及 $k$，若 $k>n$，则
\begin{align*}
\binom{n}{k}=0
\end{align*}
下列展开式在关于二项系数的讨论中很有用
\begin{align*}
(1+x)^n=\sum_{r=0}^{n}\binom{n}{r}x^r
\end{align*}
令 $x=1$ 得到
\begin{align*}
\binom{n}{0}+\binom{n}{1}+\binom{n}{2}+\cdots+\binom{n}{n}=2^n
\end{align*}
利用幂级数乘法又可以证明，对一切正整数 $a,b$，成立
\begin{align}
\binom{a}{0}\binom{b}{n}+\binom{a}{1}\binom{b}{n-1}+\cdots+\binom{a}{n}\binom{b}{0}=\binom{a+b}{n}\label{align:ch01-范德蒙德}
\end{align}
特别地
\begin{align*}
\binom{n}{0}\binom{n}{n}+\binom{n}{1}\binom{n}{n-1}+\cdots+\binom{n}{n}\binom{n}{0}=\binom{2n}{n}
\end{align*}
即
\begin{align}
\binom{n}{0}^2+\binom{n}{1}^2+\cdots+\binom{n}{n}^2=\binom{2n}{n}\label{align:ch01-二项系数平方和}
\end{align}

把排列公式推广到 $r$ 是正整数而 $n$ 是任意实数 $x$ 的场合，有时是需要的，这时记
\begin{align*}
A_x^r=x(x-1)(x-2)\cdots(x-r+1)
\end{align*}
同样定义
\begin{align*}
\binom{x}{r}=\frac{A_x^r}{r!}=\frac{x(x-1)(x-2)\cdots(x-r+1)}{r!}
\end{align*}
约定 $\binom{x}{0}=1$。

不难验算：
\begin{align}
\binom{-a}{k}=(-1)^k\binom{a+k-1}{k}\label{align:ch01-负二项系数}
\end{align}

这时，对任意实数 $\alpha$，有\textbf{Newton 二项式}
\begin{align*}
(1+x)^\alpha=\sum_{r=0}^{\infty}\binom{\alpha}{r}x^r
\end{align*}
当 $\alpha$ 为正整数 $n$ 时，它化为上述展开式。

\subsection*{三、概率直接计算的例子}

\begin{example}[文集排列]\label{example:ch01-文集排列}
一部四册的文集按任意次序放到书架上去，问各册自右向左或自左向右恰成 $1,2,3,4$ 的顺序的概率是多少？

\textbf{解}~~若以 $a,b,c,d$ 分别表示自左向右排列的书的册号，则上述文集放置的方式可与向量 $(a,b,c,d)$ 建立一一对应，因为 $a,b,c,d$ 取值于 $1,2,3,4$，因此这种向量的总数相当于 $4$ 个元素的全排列数 $4!=24$，由于文集按“任意的”次序放到书架上去，因此这 $24$ 种排列中出现任意一种的可能性都相同，这是古典概型概率，其有利场合有 $2$ 种，即自左向右或自右向左成 $1,2,3,4$ 顺序，因此所求概率为 $\frac{2}{24}=\frac{1}{12}$。
\end{example}

\begin{example}[幸运37选7]\label{example:ch01-彩票}
在幸运 $37$ 选 $7$ 福利彩票中，每期从 $1,2,\cdots,37$ 中开出 $7$ 个基本号码和一个特殊号码，彩民们在购买每一张彩票时都预先选定 $7$ 个号码。规定 $7$ 个基本号码全部选中者获一等奖，选中 $6$ 个基本号码及特殊号码者获二等奖。试求购买一张彩票中一等奖的概率 $p_1$ 及中二等奖的概率 $p_2$。

\textbf{解}~~把从 $37$ 个数中取 $7$ 个的各种取法作为样本点全体，这是古典概型；总数虽大，为 $\binom{37}{7}=10295472$，但还是有限的；由于摇奖时各数地位的对称性，因此这些样本点的出现是等可能的。

一等奖的有利场合数目为 $\binom{7}{7}\binom{30}{0}=1$，故
\begin{align*}
p_1=\frac{1}{\binom{37}{7}}=9.713\times 10^{-8}
\end{align*}
即中一等奖的概率约为一千万分之一。

二等奖的有利场合数目为 $\binom{7}{6}\binom{1}{1}\binom{29}{0}=7$，故
\begin{align*}
p_2=\frac{7}{\binom{37}{7}}=7\times p_1=6.8\times 10^{-7}
\end{align*}
\end{example}

\begin{example}[硬币比较]\label{example:ch01-硬币比较}
甲有 $n+1$ 个硬币，乙有 $n$ 个硬币，双方投掷之后进行比较，求甲掷出的正面数比乙掷出的正面数多的概率。

\textbf{解}~~这个问题初看非经繁复计算难求答案，但是若充分利用它特有的对称性并选择适当的样本空间，则能迅速求解。

若以 $A$ 记“甲的正面数$>$乙的正面数”，则 $\overline{A}$ 表示“甲的正面数$\leqslant$乙的正面数”，但是考虑到甲只比乙多一个硬币这一特殊情况，则 $\overline{A}$ 又表示“甲的反面数$>$乙的反面数”。再由硬币的对称性，显然 $P(A)=P(\overline{A})$。因此我们如果构造只有两个样本点的样本空间 $\Omega=\{A,\overline{A}\}$，便可轻而易举地写出答案 $P(A)=\frac{1}{2}$。
\end{example}

\begin{example}[投球入格]\label{example:ch01-投球入格}
设有 $n$ 个球，每个球都能以同样的概率 $\frac{1}{N}$ 落到 $N$ 个格子（$N\geqslant n$）的每一个格子中，试求：

（1）某指定的 $n$ 个格子中各有一个球的概率；

（2）任何 $n$ 个格子中各有一个球的概率。

\textbf{解}~~这是一个古典概型问题，由于每个球可落入 $N$ 个格子中的任一个，所以 $n$ 个球在 $N$ 个格子中的分配相当于从 $N$ 个元素中选取 $n$ 个进行有重复的排列，故共有 $N^n$ 种可能分配。

在第一个问题中，有利场合相当于 $n$ 个球在那指定的 $n$ 个格子中全排列，总数为 $n!$，因而所求概率为
\begin{align*}
P_1=\frac{n!}{N^n}
\end{align*}

在第二个问题中，$n$ 个格子可以任意，即可以从 $N$ 个格子中任意选出 $n$ 个来，这种选法共有 $\binom{N}{n}$ 种，对于每种选定的 $n$ 个格子，有利场合正如第一个问题一样为 $n!$，故所求概率为
\begin{align*}
P_2=\frac{\binom{N}{n}n!}{N^n}=\frac{N!}{N^n(N-n)!}
\end{align*}

这个例子是古典概型中一个很典型的问题，不少实际问题都可以归结为它。

例如，若把球解释为粒子，把格子解释为相空间中的小区域，则这个问题便相应于统计物理学中的 Maxwell--Boltzmann 统计。

这也联系着概率论历史上有名的\textbf{生日问题}：求参加某次集会的 $n$ 个人中至少有两个人生日相同的概率 $p_n$。若把 $n$ 个人看作上面问题中的 $n$ 个球，而把一年的 $365$ 天作为格子，则 $N=365$，这时所求的概率就是 $1-P_2$，即
\begin{align*}
p_n=1-365\times 364\times\cdots\times(365-n)/365^n
\end{align*}
下表给出若干 $n$ 与 $p_n$ 的数值：
\begin{center}
\begin{tabular}{c|ccccccc}
$n$ & 5 & 10 & 20 & 23 & 30 & 40 & 60 \\
\hline
$p_n$ & 0.027 & 0.117 & 0.411 & 0.507 & 0.706 & 0.891 & 0.994
\end{tabular}
\end{center}
当 $n=23$ 时，同生日的概率就超过 $\frac{1}{2}$；当 $n=40$ 时，该概率几达九成；进一步当 $n=60$ 时，几乎可以肯定必有二人生日。总之，投球问题中球相遇的概率比预料的大得多，这种意外在研究随机现象中时常遇见，也算是随机现象的特性之一吧！
\end{example}

\subsection*{四、抽签与顺序无关}

抽签是人为地引进随机性的一个最简单的例子，不单在体育比赛中广泛采用，而且在日常生活中也时常可见。关于抽签所体现的公平性，即结果虽然不同但机会却是均等的这种概念，只有通过概率论才能清楚阐明。

进一步可以说，抽签是随机化方法的一种特例，随机化方法在抽样调查、试验设计和决策中都有广泛应用，是人类可以主动利用随机性的一个例证。

下面的摸球问题值得读者仔细推敲。

\begin{example}[抽签与顺序无关]\label{example:ch01-抽签}
口袋中有 $a$ 只黑球，$b$ 只白球，它们除颜色不同外，其他方面没有差别，现在把球随机地一只只摸出来，求第 $k$ 次摸出的一只球是黑球的概率（$1\leqslant k\leqslant a+b$）。

\textbf{[第一种解法]}~~把 $a$ 只黑球及 $b$ 只白球都看作是不同的（例如设想把它们进行编号），若把摸出的球依次放在排列成一直线的 $a+b$ 个位置上，则可能的排列法相当于把 $a+b$ 个元素进行全排列，总数为 $(a+b)!$，把它们作为样本点全体。有利场合数为 $a\times(a+b-1)!$，这是因为第 $k$ 次摸得黑球有 $a$ 种取法，而另外 $(a+b-1)$ 次摸球相当于 $a+b-1$ 只球进行全排列，有 $(a+b-1)!$ 种构成法，故所求概率为
\begin{align*}
P_k=\frac{a\times(a+b-1)!}{(a+b)!}=\frac{a}{a+b}
\end{align*}
这个结果与 $k$ 无关！

细想一下，就会发觉这个结果与我们平常的生活经验是一致的。例如在体育比赛中进行抽签，对各队机会均等，与抽签的先后次序无关。

\textbf{[第二种解法]}~~把 $a$ 只黑球看作是没有区别的，把 $b$ 只白球也看作是没有区别的。仍把摸出的球依次放在排列成一直线的 $a+b$ 个位置上，因若把 $a$ 只黑球的位置固定下来则其他位置必然是放白球，而黑球的位置可以有 $\binom{a+b}{a}$ 种放法，以这种放法作为样本点。这时有利场合数为 $\binom{a+b-1}{a-1}$，这是由于第 $k$ 次摸得黑球，这个位置必须放黑球，剩下的黑球可以在 $a+b-1$ 个位置上任取 $a-1$ 个位置，因此共有 $\binom{a+b-1}{a-1}$ 种放法。故所求概率为
\begin{align*}
P_k=\frac{\binom{a+b-1}{a-1}}{\binom{a+b}{a}}=\frac{a}{a+b}
\end{align*}

两种不同的解法答案相同！

注意考察一下两种解法的不同，就会发现主要在于选取的样本空间不同。在前一种解法中把球看作是“有个性的”，而在后一种解法中则对同色球不加区别，因此在第一种解法中要顾及各黑球间及各白球间的顺序而用排列，第二种解法则不注意顺序而用组合，但最后还是得出了相同的答案。

这种情况的产生并不奇怪，这说明对于同一随机现象，可以用不同的模型来描述，只要方法正确，结论总是一致的。在这个例子中，第二种解法中的每一个样本点是由第一种解法中的 $a!\cdot b!$ 个样本点合并而成的。

这个例子还告诉我们，在计算样本点总数及有利场合数时，必须对同一个确定的样本空间考虑，因此其中一个考虑顺序，另一个也必须考虑顺序，否则结果一定不正确。

既然同一个随机现象可用不同的样本空间来描述，因此对同一个概率也常常有多种不同的求法，我们应逐步训练自己能采用最简便的方法解题，为此熟悉同一问题的多种不同解法是重要的。

例如，对例~\ref{example:ch01-抽签} 就存在着多种不同的解法，上面提供的只是比较自然的两种。注意到在这两种解法中，我们对不同的 $k$ 用的是同一个样本空间，也就是说，我们构造了一个可以描述 $a+b$ 次摸球的样本空间，并利用它一举解决了“第 $k$（$1\leqslant k\leqslant a+b$）次摸得黑球”这一概率的计算。假如允许对不同的 $k$ 用不同的样本空间，则我们完全可以构造一个只包含前 $k$ 次试验，甚至只包含第 $k$ 次试验的样本空间，这时也能求得有关概率。特别是选用最后一种样本空间并利用对称性马上可以看出正确答案，不过这种做法对初学者或许不那么容易理解。
\end{example}

\subsection*{五、二项分布与超几何分布}

产品抽样检查有两类，即有放回抽样与不放回抽样。在有放回抽样中，被抽出的产品检验后仍放回产品中，再抽第二次，因此这件产品以后仍然可能再次被抽到。更常用的是第二类方法，即不放回的抽样方法，这时被抽到的产品不再放回，因而以后不会再被抽到。与此相应地，我们的摸球模型也假定为有放回与不放回摸球两类，这两个情形得到的结果是不同的。

下面是古典概型概率计算中的一个典型问题，它有着多方面应用，特别在产品检验方面起很大作用。

\begin{example}[二项分布与超几何分布]\label{example:ch01-二项与超几何}
如果某批产品中有 $a$ 件次品 $b$ 件合格品，我们采用有放回及不放回抽样方式从中抽 $n$ 件产品，问正好有 $k$ 件是次品的概率各是多少？

所求的概率显然与抽样方式有关，下面我们分别来讨论。

\textbf{[有放回抽样场合]}~~把 $a+b$ 件产品进行编号，有放回抽 $n$ 次，把可能的重复排列全体作为样本点，总数为 $(a+b)^n$，其中有利场合（即次品正好出现 $k$ 次）的数目是 $\binom{n}{k}a^k b^{n-k}$，故所求概率为
\begin{align}
b_k=\frac{\binom{n}{k}a^k b^{n-k}}{(a+b)^n}=\binom{n}{k}\left(\frac{a}{a+b}\right)^k\left(\frac{b}{a+b}\right)^{n-k}\label{align:ch01-二项分布}
\end{align}
$b_k$ 是二项式 $\left(\frac{a}{a+b}+\frac{b}{a+b}\right)^n$ 展开式的一般项，上述概率称为\textbf{二项分布}（分布一词的意义将在第三章阐明）。关于二项分布更一般的讨论在以后各章陆续进行。

\textbf{[不放回抽样场合]}~~从 $a+b$ 件产品中取出 $n$ 件产品的可能组合全体作为样本点，总数为 $\binom{a+b}{n}$，有利场合数为 $\binom{a}{k}\binom{b}{n-k}$，故所求概率为
\begin{align}
h_k=\frac{\binom{a}{k}\binom{b}{n-k}}{\binom{a+b}{n}}\label{align:ch01-超几何分布}
\end{align}
这个概率称为\textbf{超几何分布}。

从直观上看，当产品总数很大而抽样数不大时，采用有放回抽样与采用不放回抽样，差别应该不大。

事实上，因为
\begin{align*}
h_k&=\frac{\binom{a}{k}\binom{b}{n-k}}{\binom{a+b}{n}}=\frac{\frac{A_a^k}{k!}\cdot\frac{A_b^{n-k}}{(n-k)!}}{\frac{A_{a+b}^n}{n!}}=\binom{n}{k}\frac{A_a^k A_b^{n-k}}{A_{a+b}^n}\\
&=\binom{n}{k}\frac{a^k b^{n-k}}{(a+b)^n}\cdot\frac{\frac{A_a^k}{a^k}\cdot\frac{A_b^{n-k}}{b^{n-k}}}{\frac{A_{a+b}^n}{(a+b)^n}}
\end{align*}
而当 $k$ 比 $a$ 小得多，$n-k$ 比 $b$ 小得多时：
\begin{align*}
\frac{\frac{A_a^k}{a^k}\cdot\frac{A_b^{n-k}}{b^{n-k}}}{\frac{A_{a+b}^n}{(a+b)^n}}\approx 1
\end{align*}
因此
\begin{align*}
h_k\approx b_k
\end{align*}

在实际工作中，抽样一般都采用不放回方式，因此计算次品数为 $k$ 的概率时应该用超几何分布，但~\eqref{align:ch01-超几何分布} 的数值计算较繁复。若产品数甚大而抽样数不太大，则可利用上述性质，计算二项分布作为近似值，这时有许多专门表格可查，这样可以大大节省计算工作量。

利用上例结果，马上可以计算下列概率：若一批产品共有 $N$ 件，其中有次品 $M$（$M<N$）件，今抽取 $n$ 件，则其中恰有 $m$ 件次品的概率是
\begin{align}
P_m=\frac{\binom{M}{m}\binom{N-M}{n-m}}{\binom{N}{n}},\qquad
\begin{array}{l}
0\leqslant n\leqslant N,\,0\leqslant m\leqslant M\\
0\leqslant n-m\leqslant N-M
\end{array}\label{align:ch01-超几何分布常见形式}
\end{align}
这是超几何分布的另一种常见形式。

学到这里，细心的读者可能会发觉这样一个矛盾，在我们前面讨论中都假定产品中的次品数已知，然后根据它来计算种种概率，而在实际问题中，情况恰恰相反，次品数是未知的，并且正是我们希望通过抽样检查来确定的。

这个矛盾可通过下面办法来解决。

不难理解，抽出来的样本的质量情况在某种程度上反映了整批产品的质量情况，例如，如果整批产品中次品很多，则抽查的样本中含有次品的可能性就相当大；反之，若产品中极少次品，则从中抽查一两件产品而得到次品的可能性就很小，因而样本中所含次品数的多少就为我们估计整批产品中的次品数提供了某种信息。例如为了确定某批产品的次品率，通常采用的方法是从这批产品中抽若干件产品作为样本来检验，并用样本的次品率来估计整批产品的次品率。关于这个课题的研究，构成了数理统计的重要内容。

由于抽样带有随机性，因而不同的抽样可能得到不同的结果，所以我们有必要对各种结果出现的可能性大小进行讨论，这为我们根据样本情况推断整批产品情况提供了理论依据，这种研究是概率论的任务。从这里也看出，概率论与数理统计有着很密切的联系。

\end{example}

\begin{example}[鱼群估计]\label{example:ch01-鱼群估计}
从某鱼池中捕得 $1\,200$ 条鱼，做了红色的记号之后再放回池中，经过适当的时间后，再从池中捕 $1\,000$ 条鱼，数出其中有红色记号的鱼的数目，共有 $100$ 条，试估计鱼池中共有多少条鱼？

\textbf{解}~~设池中共有 $n$ 条鱼，$n$ 未知，是我们要估计的。更一般地，设第一次捕得的鱼有 $n_1$ 条（针对本例，$n_1=1\,200$），第二次捕得 $r$ 条（针对本例，$r=1\,000$），而其中有记号的有 $k$ 条（针对本例，$k=100$）。

现在，在第二次捕鱼中有 $k$ 条有记号鱼的概率按超几何分布~\eqref{align:ch01-超几何分布常见形式} 给出：
\begin{align*}
p_k(n)=\binom{n_1}{k}\binom{n-n_1}{r-k}\bigg/\binom{n}{r}
\end{align*}
因为实际上在 $r$ 条鱼中有 $k$ 条有记号，因此我们求 $n$ 使得上式概率达到最大，并把这个数值作为池中鱼数的估计。

由于
\begin{align*}
\frac{p_k(n)}{p_k(n-1)}=\frac{(n-n_1)(n-r)}{(n-n_1-r+k)n}=\frac{n^2-nn_1-nr+n_1 r}{n^2-nn_1-nr+nk}=\rho
\end{align*}
因此当 $nk<n_1 r$ 时，$\rho>1$，而当 $nk>n_1 r$ 时，$\rho<1$，即 $p_k(n)$ 当 $n<\frac{n_1 r}{k}$ 时是 $n$ 的增函数，而当 $n>\frac{n_1 r}{k}$ 时是 $n$ 的减函数，所以当 $n$ 等于
\begin{align*}
\hat{n}=\left[\frac{n_1 r}{k}\right]
\end{align*}
时，$p_k(n)$ 达到最大值。这样我们把 $\left[\frac{n_1 r}{k}\right]$ 作为鱼池中鱼总数 $n$ 的估计量。

在这个例子中：
\begin{align*}
\hat{n}=\frac{1\,200\times 1\,000}{100}=12\,000
\end{align*}

这里解题的基本思路是把概率 $p_k(n)$ 看作未知参数 $n$ 的函数，称为\textbf{似然函数}（likelihood function），再通过求其最大值而得到 $n$ 的估计，这是数理统计中著名的\textbf{最大似然估计法}。这个方法可以溯源到 Laplace 及 Gauss。
\end{example}

\subsection*{六、概率的基本性质}

古典概型中，若样本空间为 $\Omega=\{\omega_1,\omega_2,\cdots,\omega_n\}$，分别以 $n(A),n(B),n(A\cup B)$ 及 $n(AB)$ 记事件 $A,B,A\cup B$ 及 $AB$ 所包含的样本点数，则显然成立计数公式
\begin{align*}
n(A\cup B)=n(A)+n(B)-n(AB)
\end{align*}
两边除以 $n$，即得概率的等式
\begin{align*}
P(A\cup B)=P(A)+P(B)-P(AB)
\end{align*}
以后将会看到，这是概率论中普遍成立的一个等式。

如果 $A$ 与 $B$ 不相容，即 $AB=\varnothing$，立刻得到
\begin{align}
P(A+B)=P(A)+P(B)\label{align:ch01-有限可加}
\end{align}

不难把这个结果推广到有限个事件的场合，即若 $A_1,A_2,\cdots,A_m$ 是 $m$ 个两两互不相容的事件，则
\begin{align}
P(A_1+A_2+\cdots+A_m)=P(A_1)+P(A_2)+\cdots+P(A_m)\label{align:ch01-有限可加推广}
\end{align}
这个结论可借助~\eqref{align:ch01-有限可加} 用归纳法证明。

\begin{theorem}[逆事件概率]\label{theorem:ch01-逆事件概率}
对任何事件 $A$，成立
\begin{align}
P(\overline{A})=1-P(A)\label{align:ch01-逆事件概率}
\end{align}
\end{theorem}

\begin{proof}
由于 $A+\overline{A}=\Omega$，$A\overline{A}=\varnothing$，所以
\begin{align*}
P(\Omega)=P(A+\overline{A})=P(A)+P(\overline{A})=1
\end{align*}
因此
\begin{align*}
P(\overline{A})=1-P(A)
\end{align*}
\end{proof}

对立事件概率等式~\eqref{align:ch01-逆事件概率} 很简单，意义也很明确，但相当有用，它为计算某些事件的概率提供了很大方便。

\begin{example}[De Méré 问题]\label{example:ch01-de-mere}
一颗骰子投 $4$ 次至少得到一个六点与两颗骰子投 $24$ 次至少得到一个双六，这两件事中哪一件有更多的机会遇到？

\textbf{解}~~以 $A$ 表示一颗骰子投 $4$ 次至少得到一个六点这一事件，为求 $P(A)$，在这种场合最方便的方法是利用~\eqref{align:ch01-逆事件概率} 式，先求 $P(\overline{A})$，这时 $\overline{A}$ 表示投一颗骰子 $4$ 次都没有出现六点，因此不难得出
\begin{align*}
P(\overline{A})=\frac{5^4}{6^4}
\end{align*}
从而得到
\begin{align*}
p_1=P(A)=1-\left(\frac{5}{6}\right)^4=0.517\,7
\end{align*}

若以 $B$ 表示两颗骰子投 $24$ 次至少得到一个双六这一事件，则用同样的方法可以求得
\begin{align*}
p_2=P(B)=1-\left(\frac{35}{36}\right)^{24}=0.491\,4
\end{align*}
这样一来，我们知道前者的机会大于 $\frac{1}{2}$，而后者的机会小于 $\frac{1}{2}$。

这个问题在概率论发展史上颇有名气，因为它是 De Méré 向 Pascal（$1623$—$1662$）提出的问题之一。正是这些问题导致了 Pascal 的研究和他与 Fermat（$1601$—$1665$）的著名通信。前已指出，他们的研究标志着概率论的诞生。
\end{example}

\begin{example}[换球问题]\label{example:ch01-换球}
一口袋中装有 $N-1$ 只黑球及 $1$ 只白球，每次从袋中随机地摸出一球，并换入一只黑球，这样继续下去，问第 $k$ 次摸球时摸到黑球的概率是多少？

\textbf{解}~~若以 $A$ 表示第 $k$ 次摸到黑球这一事件，则 $\overline{A}$ 表示第 $k$ 次摸到白球。现在计算 $P(\overline{A})$。

因为袋中只有一只白球，而每次摸出球总是换入黑球，故为了在第 $k$ 次摸到白球，则前面的 $k-1$ 次一定不能摸到白球。因此 $\overline{A}$ 等价于下列事件：在前 $k-1$ 次摸球时都摸出黑球而第 $k$ 次摸出白球，这一事件的概率为
\begin{align*}
\frac{(N-1)^{k-1}\cdot 1}{N^k}=\left(1-\frac{1}{N}\right)^{k-1}\cdot\frac{1}{N}
\end{align*}
这样
\begin{align*}
P(A)=1-P(\overline{A})=1-\left(1-\frac{1}{N}\right)^{k-1}\cdot\frac{1}{N}
\end{align*}

公式~\eqref{align:ch01-逆事件概率} 的应用使这个题目很快获得解决，假如直接计算 $P(A)$ 则困难得多。从公式~\eqref{align:ch01-逆事件概率} 中容易看出，只要 $P(A)$ 或 $P(\overline{A})$ 中的任何一个知道了，则可以求得另一个。在不同问题中，有的求 $P(A)$ 容易，求 $P(\overline{A})$ 困难；有的正好相反。利用式~\eqref{align:ch01-逆事件概率}，我们就可以先求容易的一个，再去求另一个。
\end{example}

从古典概型的概率研究中，我们发现概率有下面三个基本性质：

（i）对于任何事件 $A$，$P(A)\geqslant 0$；

（ii）$P(\Omega)=1$；

（iii）若 $A_1,A_2,\cdots,A_m$ 两两互不相容，则
\begin{align*}
P(A_1+A_2+\cdots+A_m)=P(A_1)+P(A_2)+\cdots+P(A_m)
\end{align*}

第一个性质称为概率的\textbf{非负性}，第二个性质称为概率的\textbf{规范性}，第三个性质称为概率的\textbf{（有限）可加性}。

\end{document}